\section{从偏好到奖励模型}
\label{sec:17-Constrained-Guarantees/05-Learning-from-Comparisons/02}

一份训练日志摆在面前，两列数字方向相反。左边一列是奖励模型给当前策略打的平均分，从 \(0.4\) 一路涨到 \(2.9\)，曲线平滑，没有任何异常。右边一列是每两百步做一次的人工抽检胜率，从 \(58\%\) 涨到 \(66\%\)，然后掉头往下，最后回到 \(52\%\)。

训练没有发散，梯度范数正常，学习率也没动过。奖励在涨，而被奖励衡量的那个东西在跌。

要弄清这是怎么发生的，得先看清奖励模型到底是从上一节那个比较模型里继承了什么。

\subsection{同一个似然，换一种参数化}

上一节给每个物品配一个自由参数 \(s_i\)。现在物品换成``在输入 \(x\) 下的一个回答 \(y\)''，这样的对有无穷多个，不可能每个配一个参数。于是把分数换成一个函数：\(r_\theta(x,y)\)。Bradley--Terry 的形式原样保留，
\[
P(y_w\succ y_l\given x)=\sigma\bigl(r_\theta(x,y_w)-r_\theta(x,y_l)\bigr),
\]
其中 \(y_w\) 是标注员选中的回答，\(y_l\) 是被弃的那个。取负对数，逐条比较的损失是
\[
\mathcal{L}(\theta)=-\log\sigma\bigl(r_\theta(x,y_w)-r_\theta(x,y_l)\bigr).
\]

这就是奖励模型的训练目标，一行不多。它不是一个新发明的损失，它就是 Bradley--Terry 的负对数似然，只不过 \(s_i\) 被换成了 \(r_\theta(x,y)\)。\hyperref[sec:13-Machine-Learning/01-Learning-Problems-and-Evaluation/03]{``目标函数与噪声模型的对应：从生成机制到优化目标''}讲的那条对应关系在这里表现得很干净：先写下``标注员怎样产生选择''这个噪声模型，损失函数就被确定了，没有别的自由度可调。

换参数化会丢掉一样东西：凸性。\(\ell\) 对 \(s\) 是凹的，但 \(r_\theta\) 对 \(\theta\) 是非线性的，复合之后目标对 \(\theta\) 不再凹。上一节那句``最大似然的渐近性质原则上都适用''，到这里连``原则上''也不剩多少。而上一节的三条限制——只有差值可辨识、比较图的连通性、传递性假设——一条都没有被换掉。它们只是变得更难看见了。

\subsection{被拟合的是标注分布，不是``好''}

模型拟合的目标写得很明确：\(P(y_w\succ y_l\given x)\) 这个条件分布。这个分布来自实际发生的标注过程，所以它把整条流水线的性质都吸收了进去。

标注指南写了什么，就定义了什么叫更好。一条``回答应尽量完整''的说明会让长回答系统性地占优，于是 \(r_\theta\) 学到的偏好里就含有一个长度项——这不是模型的缺陷，这是它在正确地拟合被给定的分布。标注员的构成同样进入分布：什么语言、什么专业背景、在什么时段工作，都会改变边际判断的走向。而比较对怎么采样决定了这个分布在哪些区域被观测到，上一节讲的比较图连通性在这里的形态是：只在同一个 \(x\) 下比较两个回答，就等于从不采集跨 \(x\) 的边。

于是有一句判断需要被明确写下来：\textbf{奖励模型学到的是标注者偏好的分布，不是``好''本身}。它与``好''之间的差距，与\hyperref[sec:13-Machine-Learning/01-Learning-Problems-and-Evaluation/02]{``损失、风险、目标与指标不是同一个量''}讲的是同一道缝隙——指标是目标的代理，把代理优化到极致不会让缝隙缩小，只会让缝隙的形状被完整地暴露出来。\hyperref[sec:18-Synthesis/01-System-Lifecycle/01]{``把诉求写成目标函数与约束''}里那条更硬的说法在这里同样成立：目标函数会被精确地钻空子，而钻的正是翻译时留下的那个空。

\subsection{绝对尺度没有定义}

上一节那个整体平移的自由度，在这里以更隐蔽的方式存在：把 \(r_\theta\) 换成 \(r_\theta+c(x)\)，其中 \(c\) 是任意只依赖输入的函数，同一个 \(x\) 下的所有差值不变，损失一模一样。所以\textbf{奖励的绝对值不由数据决定，每个输入上都各有一个未定的偏移}。

这条限制有直接的操作后果。给奖励设一个绝对阈值来筛数据（``保留 \(r>1.5\) 的样本''），比较不同 \(x\) 上的奖励值来分配训练权重，或者把奖励值当作质量分数直接汇报——这三件事都在使用一个未定义的量。它们不会报错，会给出一个数，而那个数的大小是训练时的随机初始化和正则项定出来的，换一次种子就会变。只有同一个 \(x\) 下两个回答的奖励之差才是数据约束过的量。

顺带说清一件常被混淆的事：奖励模型的输出经过 \(\sigma\) 之后可以被校准，也确实应当检查校准——若模型说两个回答的胜率是 \(0.7\)，在类似的比较对上应当真有大约 \(0.7\) 的比例被这样标注，\hyperref[sec:11-Mathematical-Statistics/07-Statistical-Decision-and-Reliable-Prediction/04]{``Proper scoring rule 与校准衡量概率是否诚实''}给出的诊断在这里可以原样使用。但被校准的是``胜率''这个成对的量，不是单个 \(r_\theta(x,y)\)。校准检查通不过说明模型有问题，通过了也不能让绝对尺度获得含义。

\subsection{优化会把策略推出奖励模型的见闻范围}

现在回到开头那份日志。

奖励模型的训练数据来自某个初始策略产生的回答。它在那片区域上见过大量比较，估得也准。随后这个奖励模型被拿去当优化目标，而优化的作用恰恰是把策略往高奖励的方向推——推向的正是奖励模型给分最高的区域，也就是它见过的样本最稀疏、外推最不可靠的区域。奖励值继续上涨，而支撑这个数值的证据在同时变薄。

这个结构在本书里出现过。\hyperref[sec:16-Control-and-Reinforcement-Learning/02-Reinforcement-Learning/08]{``Off-policy 校正用权重交换分布偏差与方差''}处理的是同一件事：用一个策略产生的数据去评价另一个策略，两个分布的支撑不重合时，重要性权重会爆，估计的方差失控。而\hyperref[sec:16-Control-and-Reinforcement-Learning/02-Reinforcement-Learning/06]{``Deadly Triad 让自举误差在分布外放大''}讲的是这种分布外误差怎样通过优化回路被逐步放大而不是被平均掉。奖励过优化是这两条的合成：一个只在旧分布上可信的函数，被一个不断改变分布的过程当作目标。\hyperref[sec:13-Machine-Learning/01-Learning-Problems-and-Evaluation/04]{``数据划分、泄漏与分布偏移''}给出的那句判断在这里最简洁——评估集必须真的模拟未见数据，而当优化本身在制造未见数据时，任何固定的评估集都会逐步失效。

常见的抑制手段是在目标里加一项对初始策略的 KL 惩罚，\hyperref[sec:07-Information-Theory/03-Divergences/01]{``错误模型的代价：从交叉熵到 KL 散度''}里那个量在这里的作用是当作一根绳子：不许策略跑出奖励模型有证据的那片区域太远。这确实有效，而它的性质要说清楚——它是一个约束，不是一个修正。它没有让奖励模型在远处变准，只是不让人走到远处去。绳子该放多长，没有理论答案，靠人工抽检来定。

\begin{center}
\begin{tikzpicture}[x=1cm,y=1cm,font=\small,>=stealth]
  \draw[->,black!60] (0,0) -- (5.9,0);
  \draw[->,black!60] (0,0) -- (0,3.3);
  \draw[blue!70,thick,domain=0:0.98,samples=70,variable=\x]
    plot ({\x*5.4},{0.6+2.2*\x*(2-\x)});
  \draw[orange!80!black,thick,dashed,domain=0:0.9,samples=70,variable=\x]
    plot ({\x*5.4},{0.6+3.6*\x-4.4*\x*\x});
  \draw[black!45,dotted] (2.21,0) -- (2.21,1.34);
  \fill[black!70] (2.21,1.34) circle (0.06);
  \node[font=\footnotesize,anchor=north] at (2.21,-0.08) {过优化开始};
  \node[font=\footnotesize,anchor=north] at (5.0,-0.08) {优化强度};
  \node[font=\footnotesize,black!75,anchor=west] at (3.2,2.85) {奖励模型分数};
  \node[font=\footnotesize,black!75,anchor=west] at (2.6,0.72) {人工抽检的真实质量};
  \node[font=\footnotesize,anchor=south east] at (-0.05,3.0) {评分};
\end{tikzpicture}

\emph{两条线在同一段区间内一起上升，随后分岔：实线继续涨，虚线掉头。分岔点的位置随任务、奖励模型规模与数据量而变，这条曲线是被反复观察到的经验规律，不是定理。}
\end{center}

图注里那句限定不是客套。分岔点在哪里没有可计算的表达式，甚至在某些设定下真实质量的下降并不明显。这与\hyperref[sec:14-Deep-Learning/09-Scaling-and-Adaptation/01]{``损失随规模下降，但只是幂律''}处理的那类结论属于同一档：形状稳定、可以据此做工程决策，但不携带``必然如此''的保证，超出观测过的范围就不该指望它继续成立。实践中的做法因此是纯经验的——留一条人工抽检的曲线，看到它掉头就停。

这条实线与虚线的分岔还有一个更长期的版本。一旦系统上线，它产生的回答会进入下一轮的标注数据，标注员看到的分布已经被上一轮优化改变过；下一版奖励模型于是在一个由自己塑造过的分布上训练。\hyperref[sec:18-Synthesis/01-System-Lifecycle/06]{``上线之后系统开始改变它所预测的世界''}讲的正是这个回路，它让``奖励模型见闻之外''这件事变成一个随时间移动的目标，而不是一次性可以标定的边界。

\subsection{往下一步}

这一节从头到尾假定了一件事：那些比较来自``标注员''这个单数主体。实际上比较来自几十上百个人，他们的判断彼此不一致，而损失函数把所有比较不加区分地加在一起。这个求和动作看上去只是记账，实际上是在做一次偏好聚合——把多个人的偏好序合成一个分数函数。

下一节要说清楚，这个动作没有中立的做法：不存在一个在所有情形下都合理的聚合规则，这一点有明确的定理。于是``众包标注取多数''不是一个技术细节，它是一次已经被做出的选择。
